%%%%%%%%%%%%%%%%%%%%%%%%%%%%%%%%%%%%%%%%%
% Masters/Doctoral Thesis 
% LaTeX Template
% Version 2.5 (27/8/17)
%
% This template was downloaded from:
% http://www.LaTeXTemplates.com
%
% Version 2.x major modifications by:
% Vel (vel@latextemplates.com)
%
% This template is based on a template by:
% Steve Gunn (http://users.ecs.soton.ac.uk/srg/softwaretools/document/templates/)
% Sunil Patel (http://www.sunilpatel.co.uk/thesis-template/)
%
% Template license:
% CC BY-NC-SA 3.0 (http://creativecommons.org/licenses/by-nc-sa/3.0/)
%
%%%%%%%%%%%%%%%%%%%%%%%%%%%%%%%%%%%%%%%%%

%-------------------------------------------------
%	PACKAGES AND OTHER DOCUMENT CONFIGURATIONS
%-------------------------------------------------

\documentclass[
11pt, % The default document font size, options: 10pt, 11pt, 12pt
%oneside, % Two side (alternating margins) for binding by default, uncomment to switch to one side
english, % ngerman for German
singlespacing, % Single line spacing, alternatives: onehalfspacing or doublespacing
%draft, % Uncomment to enable draft mode (no pictures, no links, overfull hboxes indicated)
%nolistspacing, % If the document is onehalfspacing or doublespacing, uncomment this to set spacing in lists to single
liststotoc, % Uncomment to add the list of figures/tables/etc to the table of contents
%toctotoc, % Uncomment to add the main table of contents to the table of contents
%parskip, % Uncomment to add space between paragraphs
%nohyperref, % Uncomment to not load the hyperref package
headsepline, % Uncomment to get a line under the header
%chapterinoneline, % Uncomment to place the chapter title next to the number on one line
%consistentlayout, % Uncomment to change the layout of the declaration, abstract and acknowledgements pages to match the default layout
]{formatting} % The class file specifying the document structure

\usepackage[utf8]{inputenc} % Required for inputting international characters
\usepackage[T1]{fontenc} % Output font encoding for international characters
\usepackage{mathpazo} % Use the Palatino font by default
\usepackage[ddmmyyyy]{datetime}
\usepackage[backend=biber,style=ieee]{biblatex} % Use bibtext and IEEE reference style
\usepackage[autostyle=true]{csquotes} % Required to generate language-dependent quotes in the bibliography

\addbibresource{literatuurlijst.bib} % The filename of the bibliography

%----------------------------
%	MARGIN SETTINGS
%----------------------------

\geometry{
	paper=a4paper, % Change to letterpaper for US letter
	inner=2.5cm, % Inner margin
	outer=3.8cm, % Outer margin
	bindingoffset=.5cm, % Binding offset
	top=1.5cm, % Top margin
	bottom=1.5cm, % Bottom margin
	%showframe, % Uncomment to show how the type block is set on the page
}

%----------------------------
%	THESIS INFORMATION
%----------------------------

\thesistitle{Schaalbaarheid Developers.nl} % Your thesis title, this is used in the title and abstract, print it elsewhere with \ttitle
\technicalSupervisor{Tanja Ubert} % Your supervisor's name, this is used in the title page, print it elsewhere with \techsupname
\skillsSupervisor{Judith Lemmens} % skillssupname

\examiner{} % Your examiner's name, this is not currently used anywhere in the template, print it elsewhere with \examname
\degree{Informatica} % Your degree name, this is used in the title page and abstract, print it elsewhere with \degreename
\author{Kaj de Munter} % Your name, this is used in the title page and abstract, print it elsewhere with \authorname
\version{v0.2} % Version number, used in omslag and titelblad, print it with \versionnr

\addresses{} % Your address, this is not currently used anywhere in the template, print it elsewhere with \addressname

\subject{} % Your subject area, this is not currently used anywhere in the template, print it elsewhere with \subjectname
\keywords{} % Keywords for your thesis, this is not currently used anywhere in the template, print it elsewhere with \keywordnames
\university{Hogeschool Rotterdam} % Your university's name and URL, this is used in the title page and abstract, print it elsewhere with \univname
\department{Communicatie, Media, en Informatietechnologie} % Your department's name and URL, this is used in the title page and abstract, print it elsewhere with \deptname
\group{} % Your research group's name and URL, this is used in the title page, print it elsewhere with \groupname
\faculty{} % Your faculty's name and URL, this is used in the title page and abstract, print it elsewhere with \facname

\AtBeginDocument{
\hypersetup{pdftitle=\ttitle} % Set the PDF's title to your title
\hypersetup{pdfauthor=\authorname} % Set the PDF's author to your name
\hypersetup{pdfkeywords=\keywordnames} % Set the PDF's keywords to your keywords
}

\hoofdvraag{Op welke wijze kan Developers.nl de architectuur van hun huidige systemen beter schaal- en onderhoudbaar maken?}
\deelvraaghuidig{Hoe goed is de schaalbaarheid van de huidige websites?}
\deelvraagtechnieken{Wat voor technieken zijn er voor schaal- en onderhoudbaarheid?}
\deelvraagverbetering{Wat voor verbeteringen ten aanzien van schaal- en onderhoudbaarheid kunnen worden toegepast op de huidige websites?}
\deelvraagimplementatie{Hoe gaan deze verbeteringen geïmplementeerd worden?}
\deelvraagrequirements{Voldoen de verbeteringen aan de vereiste requirements?}
\begin{document}

\renewcommand{\tablename}{Tabel}

\frontmatter % Use roman page numbering style (i, ii, iii, iv...) for the pre-content pages

\pagestyle{plain} % Default to the plain heading style until the thesis style is called for the body content

%----------------------------
%	OMSLAG
%----------------------------

\begin{titlepage}
	\begin{center}
		
		\vspace*{.06\textheight}
		{\scshape\LARGE \univname\par}\vspace{1.5cm} % University name
		\textsc{\Large Bachelor scriptie}\\[0.5cm] % Thesis type
		
		\HRule \\[0.4cm] % Horizontal line
		{\huge \bfseries \ttitle\par}\vspace{0.4cm} % Thesis title
		\HRule \\[1.5cm] % Horizontal line
		
		\begin{minipage}[t]{0.4\textwidth}
			\begin{flushleft} \large
				\emph{Auteur:}\\\authorname\\0911825
				
			\end{flushleft}
		\end{minipage}
		\begin{minipage}[t]{0.4\textwidth}
			\begin{flushright} \large
				\emph{Begeleiders:} \\
				\techsupname\\
				\skillssupname
			\end{flushright}
		\end{minipage}\\[3cm]
		
		{\large \today}\\[1cm] % Date
		{\Large \texttt{\versionnr}}
		
		\vfill
		
		\includegraphics{HRO} % University/department logo
		
		\vfill
	\end{center}
\end{titlepage}

%----------------------------
%	TITELBLAD
%----------------------------

\begin{titlepage}
	\begin{center}
		{\scshape\LARGE Developers.nl\par}\vspace{0.5cm}
		{\scshape\LARGE \univname\par}\vspace{1.5cm} % University name
		\textsc{\Large Bachelor scriptie}\\[0.5cm] % Thesis type
		
		\HRule \\[0.4cm] % Horizontal line
		{\huge \bfseries \ttitle\par}\vspace{0.4cm} % Thesis title
		\HRule \\[1.5cm] % Horizontal line
		
		\begin{minipage}[t]{0.4\textwidth}
			\begin{flushleft} \large
				\emph{Auteur:}\\\textbf{\authorname}\\
				0911825@hr.nl\\
				k.demunter@developers.nl\\
				06-81019142\\
			\end{flushleft}
		\end{minipage}
		\begin{minipage}[t]{0.4\textwidth}
			\begin{flushright} \large
				\emph{Schoolbegeleiders:}\\
				\textbf{\techsupname}\\
				t.ubert@hr.nl\\
				\textbf{\skillssupname}\\
				j.h.i.lemmens@hr.nl\\[0.3cm]
				
				\emph{Bedrijfsbegeleiders:}\\
				\textbf{Maarten de Boer}\\
				m.deboer@developers.nl\\
				\textbf{Jelle van de Haterd}\\
				j.vandehaterd@developers.nl
			\end{flushright}
		\end{minipage}\\[2cm]
		
		\vfill
		
		\large \textit{Een scriptie ingediend ter voldoening aan de\\ vereiste competenties voor de opleiding \degreename}\\[0.3cm]
		
		\deptname % Research group name and department name
		
		\vfill
		
		{\large \today}\\[0.5cm] % Date
		{\Large \texttt{\versionnr}}\\[1cm]
		
		\includegraphics[height=2cm]{Developers} \hspace{3cm}
		\includegraphics[height=2cm]{HRO}

		\vfill
	\end{center}
\end{titlepage}

%----------------------------
%	ABSTRACT PAGE
%----------------------------

\begin{abstract}
\addchaptertocentry{\abstractname} % Add the abstract to the table of contents
Samenvatting van de scriptie die ik als laatst pas ga schrijven\ldots Lorem ipsum dolor sit amet, consetetur sadipscing elitr, sed diam nonumy eirmod tempor invidunt ut labore et dolore magna aliquyam erat, sed diam voluptua. At vero eos et accusam et justo duo dolores et ea rebum. Stet clita kasd gubergren, no sea takimata sanctus est Lorem ipsum dolor sit amet. Lorem ipsum dolor sit amet, consetetur sadipscing elitr, sed diam nonumy eirmod tempor invidunt ut labore et dolore magna aliquyam erat, sed diam voluptua. At vero eos et accusam et justo duo dolores et ea rebum. Stet clita kasd gubergren, no sea takimata sanctus est Lorem ipsum dolor sit amet. Lorem ipsum dolor sit amet, consetetur sadipscing elitr, sed diam nonumy eirmod tempor invidunt ut labore et dolore magna aliquyam erat, sed diam voluptua. At vero eos et accusam et justo duo dolores et ea rebum. Stet clita kasd gubergren, no sea takimata sanctus est Lorem ipsum dolor sit amet.
\end{abstract}

\begin{voorwoord}
\addchaptertocentry{\voorwoordname}
Bedankt aan iedereen die mij koffie heeft gebracht afgelopen periode\ldots Lorem ipsum dolor sit amet, consetetur sadipscing elitr, sed diam nonumy eirmod tempor invidunt ut labore et dolore magna aliquyam erat, sed diam voluptua. At vero eos et accusam et justo duo dolores et ea rebum. Stet clita kasd gubergren, no sea takimata sanctus est Lorem ipsum dolor sit amet. Lorem ipsum dolor sit amet, consetetur sadipscing elitr, sed diam nonumy eirmod tempor invidunt ut labore et dolore magna aliquyam erat, sed diam voluptua. At vero eos et accusam et justo duo dolores et ea rebum. Stet clita kasd gubergren, no sea takimata sanctus est Lorem ipsum dolor sit amet. 

Lorem ipsum dolor sit amet, consetetur sadipscing elitr, sed diam nonumy eirmod tempor invidunt ut labore et dolore magna aliquyam erat, sed diam voluptua. At vero eos et accusam et justo duo dolores et ea rebum. Stet clita kasd gubergren, no sea takimata sanctus est Lorem ipsum dolor sit amet.
\end{voorwoord}

%------------------------------------------------
%	LIST OF CONTENTS/FIGURES/TABLES PAGES
%------------------------------------------------

\renewcommand{\contentsname}{Inhoud}
\tableofcontents % Prints the main table of contents

\renewcommand{\listfigurename}{Figurenlijst}
\listoffigures % Prints the list of figures

\renewcommand{\listtablename}{Tabellenlijst}
\listoftables % Prints the list of tables

%----------------------------
%	GLOSSARY
%----------------------------
\begin{begrippen}{p{3cm}|p{10cm}} % Include a list of abbreviations (a table of two columns)
	\textbf{Kubernetes} & Kubernetes is an open-source system for automating deployment, scaling, and management of containerized applications \parencite{Kubernetes}.\\
	
	\textbf{Docker} & Docker is a tool designed to make it easier to create, deploy, and run applications by using containers \parencite{Docker}.\\
	
	\textbf{Ansible} & Ansible is a radically simple IT automation engine that automates cloud provisioning, configuration management, application deployment, intra-service orchestration, and many other IT needs \parencite{Ansible}.\\
	
	\textbf{Docker Swarm} & Docker Swarm provides native clustering functionality for Docker containers, which lets you turn a group of Docker engines into a single, virtual Docker engine \parencite{DockerSwarm}.\\
	
	\textbf{Infrastructure as Code} & Infrastructure as code describes the idea of using a high-level programming language to
	control IT systems \parencite{IaC}.\\
	
	\textbf{Ansible Tower} & Red Hat Ansible Tower helps you scale IT automation, manage complex deployments and speed productivity. Centralize and control your IT infrastructure with a visual dashboard, role-based access control, job scheduling, integrated notifications and graphical inventory management \parencite{AnsibleTower}.\\
	
	\textbf{Chef} & Deploy new code faster and more efficiently. Automate infrastructure and applications \parencite{Chef}.\\
	
	\textbf{Puppet} & Powerful infrastructure automation and delivery \parencite{Puppet}.\\
	
	\textbf{Terraform} &
	Terraform is a tool for building, changing, and versioning infrastructure safely and efficiently. Terraform can manage existing and popular service providers as well as custom in-house solutions \parencite{Terraform}.\\
	
	\textbf{Cloud computing providers} & Cloud computing is a model for enabling ubiquitous, convenient, on-demand network
	access to a shared pool of configurable computing resources (e.g., networks, servers,
	storage, applications, and services) that can be rapidly provisioned and released with
	minimal management effort or service provider interaction \parencite{CloudComputing}.
	
	Er bestaan veel cloud computing providers, waaronder bijvoorbeeld:
	\begin{itemize}
		\item Amazon web services
		\item DigitalOcean
		\item Microsoft Azure
		\item Google Cloud Platform
	\end{itemize}\\
	
	\textbf{12-factor app} & Een methodologie voor het bouwen van Software as a Service (SaaS) applications \parencite{12Factor}. \\
	
	\textbf{The Open Group Architecture Framework} & A generic framework to build different IT architectures frameworks \parencite{TOGAF}.\\
	
	\textbf{4+1 architectural view model} & A model for describing the architecture of software-intensive systems, based on the use of multiple, concurrent views \parencite{4plus1}. \\
\end{begrippen}


%----------------------------------------------------------------------------------------
%	ABBREVIATIONS
%----------------------------------------------------------------------------------------

\begin{abbreviations}{ll} % Include a list of abbreviations (a table of two columns)

\textbf{MI} & \textbf{M}aintainability \textbf{I}ndex\\
\textbf{CI} & \textbf{C}ontinuous \textbf{I}ntegration\\
\textbf{CD} & \textbf{C}ontinuous \textbf{D}eployment\\
\textbf{EMS} & \textbf{E}mployee \textbf{M}anagement \textbf{S}ystem\\
\textbf{IaC} & \textbf{I}nfrastructure \textbf{A}s \textbf{C}ode\\
\textbf{K8s} & \textbf{K}ubernete\textbf{s}\\
\textbf{SaaS} & \textbf{S}oftware \textbf{a}s \textbf{a} \textbf{S}ervice\\
\textbf{TOGAF} & \textbf{T}he \textbf{O}pen \textbf{G}roup \textbf{A}rchitecture \textbf{F}ramework

\end{abbreviations}

%----------------------------------------------------------------------------------------
%	THESIS CONTENT - CHAPTERS
%----------------------------------------------------------------------------------------

\mainmatter % Begin numeric (1,2,3...) page numbering

\pagestyle{thesis} % Return the page headers back to the "thesis" style

% Include the chapters of the thesis as separate files from the Chapters folder
% Uncomment the lines as you write the chapters

\renewcommand{\chaptername}{Hoofdstuk}

\chapter{Inleiding}

\label{Chapter1}


\section{Aanleiding}

\enquote{Een visitekaartje voor het bedrijf}. Dat is het uitgangspunt wanneer interne software bij Developers.nl wordt ontwikkeld. Niet alleen qua uiterlijk, maar ook van binnen moet de code, de infrastructuur en de werkmethode in topconditie zijn. Dit heeft te maken met het feit dat de code uit de website en infrastructuur van Developers.nl open-source wordt gemaakt gedurende deze stage.

Dit betekent dat er veel onderzoek moet worden uitgevoerd op de kwaliteit van de huidige website. Met name de infrastructuur, het \enquote{Continuous Integration/Continuous Deployment} (CI/CD) en het horizontaal en verticaal schalen hebben nog veel ruimte voor verbetering. Bovendien zijn er meerdere interne systemen dan alleen de website. Het onderhouden van (binnenkort) 5 systemen vereist veel tijd en moeite, dus moet dit makkelijker worden.

\section{Belang}

Het open-source maken van de website betekent dat elke potentiële klant en/of nieuwe medewerker de mogelijkheid heeft om te bekijken wat Developers.nl qua kennis in huis heeft. Het is dus van hoog belang dat de kwaliteit gewaarborgd wordt, en dat er zo veel mogelijk nieuwe en opkomende technieken worden gebruikt.

Bovendien heeft Wheeler \parencite{WhyOpenSource} geconcludeerd dat open-source software voordelen heeft als:
\begin{itemize}
	\item Betere beveiliging
	\item Betere betrouwbaarheid 
	\item Betere prestaties
	\item Betere schaalbaarheid
	\item Mindere onderhoudskosten
\end{itemize}
Ook al is de website op eerste gezicht van de buitenkant vrij eenvoudig, er moet (om indruk te wekken op potentiële klanten en nieuwe medewerkers) onder water een stevige applicatie draaien dat eigenlijk iets \enquote{te complex} is. Kortom, er is nog veel te verbeteren en toe te voegen.

Het onderhouden van (binnenkort) 5 interne systemen betekent dat er veel werk dubbel wordt uitgevoerd. Dit kost geld, en kan beter.

\pagebreak

\section{Opdrachtgever}

Deze scriptie is geschreven in opdracht van:\\[0.3cm]
\textbf{Developers.nl}\\
Stadionweg 57C\\
3077AS Rotterdam\\
info@developers.nl\\
010-3035929\\

\subsection{Core business}
Developers.nl neemt software ontwikkelaars in dienst. De ontwikkelaars die worden aangenomen zullen voornamelijk gespecialiseerd zijn in PHP, Python, Java of front-end. Ze worden uitgezet naar een klant (een extern bedrijf) die naar een ontwikkelaar zoekt. Developers.nl kiest hier voor de beste ontwikkelaar voor de taak en zal deze inzetten bij een klant. De opdrachten van de ontwikkelaars zijn op locatie van de klant en duren voornamelijk langer dan een jaar, maar op uitzondering zijn er ook kortere opdrachten. Zodra de ontwikkelaar klaar is met zijn of haar taak zal Developers.nl zo snel mogelijk een nieuwe opdracht toewijzen \parencite{Stageplan}. Concreet zegt het positioneringsprofiel \parencite{Positioneringsprofiel}: \enquote{Detachering van developers die software applicaties bouwen voor verschillende klanten.}

\subsection{Geschiedenis}
Mark Smit heeft altijd al in de detacheringshoek gewerkt. Hieruit is de inspiratie begonnen om een detacheringsbedrijf voor zichzelf op te richten in 2012. Het bedrijf heette \enquote{JC Capacity} en was onderdeel van \enquote{JC Groep}. Mark Smit was mede eigenaar met 2 anderen binnen JC Groep. Er was op dit moment alleen nog een kleine groep PHP ontwikkelaars, totdat in 2013 een aantal Java en Python ontwikkelaars zijn aangenomen. In februari 2014 is Mark Smit samen met een andere eigenaar binnen JC Group als JC Capacity afgesplitst. Dit gebeurde in verband met meningsverschillen over welke kant het bedrijf op moest doordat JC Groep meer op de zorg gefocust was. Gelijk hierna is een naamsverandering plaatsgevonden en heette het bedrijf \enquote{Developers.nl}.

Kort na de naamsverandering is het bedrijf exponentieel gegroeid, de teams PHP, Java en Python zijn individueel gegroeid, maar ook is er een nieuwe Front-End tak bijgekomen.  vanaf 2015 is heeft de andere eigenaar het bedrijf verlaten en was Mark Smit de enige eigenaar van Developers.nl De naamsverandering en de splitsing van JC Groep heeft ervoor gezorgd dat Developers.nl meer vrijheid had om veranderingen aan te brengen en naamsbekendheid op te krikken. 

In 2016 heeft Mark Smit een tweede bedrijf opgericht met de naam \enquote{Gemvision}. Dit bedrijf vereistte meer van zijn tijd op dan hij had verwacht, en heeft hierdoor vanaf 2018 zijn management taken binnen Developers.nl overgegeven aan Maarten de Boer \parencite{Stageplan}.

\subsection{Werkomgeving}
\subsubsection{Cultuur}
Developers.nl is een jonge en ambitieuze club mensen met een sterke onderlinge samenhang. Developers is een platte organisatie. Er heerst een sfeer van vrijheid, verantwoordelijkheid en \enquote{work hard, play hard}. De omgang met elkaar is informeel en persoonlijk, met het gevoel van een vriendenclub. De organisatie wordt gerund vanuit het besef dat plezier in het werk belangrijk is. De dresscode is business casual zodat klanten hetzelfde gevoel hebben van vrijheid en verantwoordelijkheid bij de ontwikkelaars. Op het kantoor luncht iedereen gezamenlijk en wordt er regelmatig geborreld na werktijd. Drinken is inbegrepen en lunch moet zelf meegebracht worden. Volgens het Competing Values Framework van Robert Quinn en Rohrbaug (1983) valt de cultuur onder mensgericht, de medewerkers staan centraal \parencite{Bedrijfscultuur}.

Er is veel aandacht voor de medewerker als mens en zijn of haar persoonlijke leven, ook buiten werktijd. De open cultuur biedt ruimte voor het inbrengen van ambities, ideeën, vragen en wensen, op welk gebied dan ook. Vanuit het principe van wederkerigheid spant de organisatie zich in om de medewerker zoveel mogelijk te helpen en een goede werkgever te zijn, zodat medewerkers zich inspannen voor hun werkgever, collega’s en klanten.

Ontwikkeling staat in alle opzichten centraal en strekt zich uit in twee richtingen: De ontwikkeling van medewerkers, zowel persoonlijk als professioneel, en de ontwikkeling van individuele en collectieve vakinhoudelijke kennis.

Developers.nl faciliteert kennis aan haar eigen ontwikkelaars in de vorm van onder andere \enquote{TechNights}, kennissessies en gezamenlijke deelname aan conferenties in het buitenland. Het zijn van een kennisorganisatie werkt twee kanten op: klanten met uitdagende opdrachten weten Developers.nl hiervoor te vinden en dat maakt de werksituatie compleet om medewerkers blijvend te boeien.

Klanten ervaren Developers.nl als oprecht geïnteresseerd, persoonlijk, analytisch, flexibel en vlot schakelend. Medewerkers zijn eerlijk en open, communicatief sterk, professioneel, betrouwbaar en intelligent. Het feit dat ze kunnen terugvallen op de collectieve kennis van de groep en dat ze bereid zijn om ook met de klant deze kennis te delen maakt ze extra waardevol. Fieldmanagers zijn erg sterk in het doorgeven van de opdracht naar uitvoerenden \parencite{Stageplan}. 

\subsubsection{Eigen omgeving}
Developers.nl heeft rond de 60 software ontwikkelaars. Deze zijn voornamelijk op een externe opdracht bij ene klant. Elke vrijdag zullen 5 van de zogenaamde \enquote{kennisambassadeurs} op kantoor zijn. Dit zijn de meest senior ontwikkelaars per team. Deze zijn dan in staat om stagiairs en/of andere medewerkers persoonlijk te helpen, maar ze zijn altijd bereikbaar via Slack of telefonisch.

De bedrijfsbegeleider voor deze stage is \textbf{Maarten de Boer}. Dit is de algemene directeur van Developers.nl en is in 2003 afgestudeerd aan de hogeschool Inholland met Strategic marketing. Aangezien Maarten zelf geen technische kennis heeft is er ook een technische begeleider aangewezen: \textbf{Jelle van de Haterd}. Jelle is senior developer, DevOps engineer en kennisambassadeur bij Developers.nl. Hij is in 2006 afgestudeerd op de Hogeschool Rotterdam met als opleiding Grafimediatechnologie \parencite{Afstudeervoorstel}.

Binnen developers.nl word iedere stagiair behandeld als een daadwerkelijk sterke medewerker die een echte bijdrage aan het bedrijf levert. Het team staat altijd voor hen klaar. Op het moment dat iemand aangeeft dat hij/zij iets nodig heeft (denk aan bijvoorbeeld een macbook voor het programmeren) zal het team er zo veel mogelijk aan doen dat er een goede oplossing wordt aangeboden. Het team staat open voor feedback en behandelen dit op serieuze wijze.

Voor technische vragen zijn de kennisambassadeurs altijd bereikbaar, bij een vraag zullen zij hun best doen om het zo duidelijk mogelijk over te brengen, en zorgen ervoor dat ik het ook daadwerkelijk begrijp. De kennisambassadeurs kijken code inhoudelijk na en geven hier feedback op \parencite{Stageplan}.

\subsection{Taken}
Tijdens de stageperiode zal de stagiair een leidende rol aannemen in een team van 2 part-time studenten en een stagiair. Er gaat onderzoek worden uitgevoerd naar de schaal- en onderhoudbaarheid van de infrastructuur aangaande de interne systemen die Developers.nl beheert. Dit omtrent onderzoek over de huidige situatie, technieken die gebruikt kunnen worden, en het automatiseren hiervan. Door het gehele proces moet er extra rekening gehouden worden met de veiligheid van de gekozen oplossing in verband met de gevoeligheid van een aantal systemen, zoals bijvoorbeeld het Employee Management System (EMS). Uit dit onderzoek zullen een aantal verbeteringen daadwerkelijk worden geïmplementeerd en uiteindelijk een advies worden uitgebracht over eventuele vorderingen.

\section{Doelstelling}
Dit onderzoek heeft tot doel het verkrijgen van kennis en inzicht over de schaal- en onderhoudbaarheid van de interne systemen die Developers.nl beheert om vervolgens deze twee factoren in de praktijk te verbeteren.

\section{Probleemstelling}
Bij Developers.nl werken veel verschillende ontwikkelaars voor een erg variabele tijd aan de interne projecten. Dit heeft te maken met het feit dat de ontwikkelaars die eraan werken vaak tussen twee opdrachten in zitten. Dit betekent dat het van hoog belang is dat een ontwikkelaar de omgevingen sneller kan opzetten en op korte termijn een kwalitatieve toevoeging kan leveren die in productie staat. Developers.nl als product owner wilt een nieuwe toevoeging of zelfs volledig nieuwe applicatie kunnen bedenken en deze zo snel mogelijk geïmplementeerd zien. Momenteel is de workflow qua integration en deployment nog niet genoeg geoptimaliseerd om ontwikkelaars deze snelheid te beloven.

\section{Requirements}
De kwaliteit van nieuwe toevoegingen moet gewaarborgd worden. Ze moeten door Developers.nl als product owner gevalideerd kunnen worden, en daarbovenop moeten deze ook automatisch getest worden op kwaliteit. Om ze te valideren moeten er voor individuele toevoegingen naast elkaar verschillende omgevingen kunnen worden opgezet die te bekijken zijn door de product owners. Dit betekent dus dat de applicatie schaalbaar moet zijn. Om schaalbaarheid te realiseren moet de onderhoudbaarheid van de infrastructuur ook op een voldoende niveau zijn, hier moeten kwaliteitsstandaarden voor onderzocht worden. De oplossing moet zo generiek mogelijk zijn en dus voor meerdere applicaties te gebruiken zijn, en dus ook bestemd zijn op applicaties met meer verkeer.

Concreet moet de oplossing;
\begin{itemize}
	\item Meerdere instanties van de applicatie naast elkaar kunnen draaien.
	\item Generiek genoeg zijn zodat meerdere applicaties er gebruik van kunnen maken.
	\item Een manier bevatten om kwaliteit automatisch te waarborgen.
	\item Automatisch kunnen schalen bij een grote hoeveelheid verkeer.
\end{itemize}

\section{Hoofd- en Deelvragen}
\subsubsection{Hoofdvraag}
\hoofdvraagname

\subsubsection{Deelvragen}
\begin{itemize}
	\item \deeltechnieken
	\item \deelhuidig
	\item \deelverbetering
	\item \deelimplementatie
	\item \deelrequirements
\end{itemize}

\section{Methodologie}
% TODO Methodologie
\textcolor{red}{TODO: Methodologie}
\begin{itemize}
	\item het type onderzoek
	\item de wijze waarop je informatie hebt vergaard
	\item de kwaliteitseisen van je bedrijf en opdrachtgever
	\item de wijze waarop je jouw bevindingen valideert 
	\item de wijze waarop je de geldigheid en betrouwbaarheid van je bronnen hebt gecontroleerd
	\item een onderbouwde keuze voor je projectmethode
\end{itemize}

De deelvraag \enquote{\deeltechnieken} Is beantwoord door......

De deelvraag \enquote{\deelhuidig} Is beantwoord door......

De deelvraag \enquote{\deelverbetering} Is beantwoord door......

De deelvraag \enquote{\deelimplementatie} Is beantwoord door......

De deelvraag \enquote{\deelrequirements} Is beantwoord door......



\section{Planning}
\begin{table}[h]
	\caption{Planning}
	\label{tab:planning}
	\centering
	\begin{tabular}{c p{12cm}}
		\toprule
		\textbf{Week} & \textbf{Taak}\\
		\midrule
			1, 2 & Skelet opzet scriptie, inleiding \\
			3, 4 & Theoretisch kader, afbakening \\
			5, 6 & \deeltechnieken \\
			7, 8 & \deelhuidig \\
			9, 10 & \deelverbetering \\
			11, 12 & \deelimplementatie \\
			13 -- 17 & Praktijk implementatie \\
			18 -- 20 & Conclusie en \deelrequirements \\
		\bottomrule\\
	\end{tabular}
\end{table}
\section{Leeswijzer}
% TODO: Leeswijzer

Lorem ipsum dolor sit amet, consectetur adipiscing elit. Aliquam ultricies lacinia euismod. Nam tempus risus in dolor rhoncus in interdum enim tincidunt. Donec vel nunc neque. In condimentum ullamcorper quam non consequat. Fusce sagittis tempor feugiat. Fusce magna erat, molestie eu convallis ut.
\chapter{Theoretisch Kader}
\label{Chapter2}

% TODO: Inleiding hoofdstuk
In dit hoofdstuk gaan we Lorem ipsum dolor sit amet, consetetur sadipscing elitr, sed diam nonumy eirmod tempor invidunt ut labore et dolore magna aliquyam erat, sed diam voluptua. At vero eos et accusam et justo duo dolores et ea rebum. Stet clita kasd gubergren, no sea takimata sanctus est Lorem ipsum dolor sit amet.

\section{Schaalbaarheid}

Mark D. Hill heeft in 1990 onderzoek gedaan naar een concrete definitie naar schaalbaarheid \parencite{WhatIsScalability}. In zijn onderzoek concludeert hij het volgende:
\begin{quote}
	\textit{
		I examined aspects of scalability, but did not find a useful, rigorous definition of it. Without such a definition, I assert that calling a system ‘scalable’ is about as useful as calling it ‘modern’. I encourage the technical community to either rigorously define scalability or stop using it to describe systems.
	}
\end{quote}
Ondertussen zijn er meerdere pogingen gedaan om schaalbaarheid te definiëren. zo zijn L. Duboc, D. S. Rosenblum en T. Wicks op deze conclusie ingegaan en hebben een poging gedaan om een framework te creëren voor karakterisering en analyse van software schaalbaarheid \parencite{ScalabilityFramework}. Zij definiëren schaalbaarheid als:
\begin{quote}
	\textit{quality of software systems characterized by the causal impact that scaling aspects of the system environment and design have on certain measured system qualities as these aspects are varied over expected operational ranges}
\end{quote}

% TODO: uitleggen scalability framework

In een onderzoek over de kenmerken van schaalbaarheid en de impact op prestatie heeft A. B. Bondy \parencite{ScalabilityCharacteristics} schaalbaarheid verdeeld in een aantal verschillende aspecten, waaronder:
\begin{itemize}
	\item \textbf{Structural scalability} (het vermogen van een systeem om uit te breiden in een gekozen dimensie zonder ingrijpende wijzigingen in de architectuur)
	\item \textbf{Load scalability} (het vermogen van een systeem om elegant te presteren naarmate het aangeboden verkeer toeneemt)
	\item \textbf{Space scalability} (het geheugenvereiste groeit niet naar \enquote{ondraaglijke niveaus} naarmate het aantal items toeneemt)
	\item \textbf{Space-time scalability} (het systeem blijft naar verwachtingen functioneren naarmate het aantal items dat het omvat toeneemt)
\end{itemize}
Bondy definieert schaalbaarheid als het vermogen van een systeem om een toenemend aantal elementen, objecten en werk gracieus te verwerken en / of vatbaar te zijn voor uitbreiding.

H. El-Rewini en M. Abd-El-Barr noemen in het boek Advanced computer architecture and parallel processing \parencite{AdvancedArchitecture} ook een aantal \enquote{onconventionele} definities:
\begin{itemize}
	\item \textbf{Size scalability} (Meet de maximale hoeveelheid processors dat een systeem kan accommoderen)
	\item \textbf{Application scalability} (de mogelijkheid om applicatiesoftware te draaien met verbeterde prestaties op een opgeschaalde versie van het systeem)
	\item \textbf{Generation scalability} (de mogelijkheid om op te schalen door het gebruik van de volgende generatie (snelle) componenten)
	\item \textbf{Heterogeneous scalability} (het vermogen van een systeem om op te schalen met behulp van hardware- en softwarecomponenten die door verschillende leveranciers zijn gemaakt)
\end{itemize}

C.B. Weinstock en J. B. Goodenough hebben een algemeen onderzoek uitgevoerd naar schaalbaarheid \parencite{OnSystemScalability}. Zij noemen in hun conclusie dat er voornamelijk twee betekenissen van het woord schaalbaarheid zijn:
\begin{enumerate}
	\item De mogelijkheid om met verhoogde werkdruk om te gaan (zonder extra resources aan een systeem toe te voegen).
	\item De mogelijkheid om met verhoogde werkdruk om te gaan door herhaaldelijk een kosteneffectieve strategie toe te passen om de mogelijkheden van een systeem uit te breiden.
\end{enumerate}
Het valt op dat een concrete definitie van schaalbaarheid alleen duidelijk te definiëren is wanneer het in meerdere verschillende soorten is opgesplitst. Daarom zal in dit onderzoek vanaf dit punt altijd worden gespecificeerd welke soort schaalbaarheid het betreft. In dit onderzoek wordt vooral de focus gelegd op de \enquote{structural scalability} en \enquote{load scalability} uit \parencite{ScalabilityCharacteristics}. \enquote{application scalability} uit \parencite{AdvancedArchitecture} heeft veel overlapping met de load scalability. Omdat load scalability iets generieker is en de twee definities van Weinstock en Goodenough \parencite{OnSystemScalability} omvat wordt deze geprefereerd boven application scalability. De overgebleven definities zijn minder relevant voor dit onderzoek aangezien ze te maken hebben met hardware, of niet volledig toepasselijk is op infrastructuur.

\subsection{Schaalbaarheid controleren}
In het onderzoek van Weinstock en Goodenough \parencite{OnSystemScalability} zijn een aantal manieren genoemd om de schaalbaarheid van een systeem te controleren.

\subsubsection{Knelpunten}
Deze knelpunten hebben vooral te maken de eerste betekenis van Weinstock en Goodenough. Er moet gecontroleerd worden op de toenemende administratieve werkdruk, de \enquote{hard-coded} limieten op capaciteit, de user interface en de complexiteitsgraad van algoritmen.

\subsubsection{Onthullen van schaalbaarheids-aannames}
Dit gaat over het onderzoeken hoe de uitbreiding van een systeem nieuwe problemen kan onthullen. Zodra een systeem zich uitbreidt is er een grotere kans op errors in de systeemconfiguratie, \enquote{zeldzame} errors komen vaker voor, is het belangrijk dat een probleem in het systeem gelokaliseerd blijft, en kan het een stuk complexer en lastiger worden om het systeem te begrijpen.

\subsubsection{Schaalstrategieën}
Hier wordt bekeken wat voor verschillende tekortkomingen een schaalstrategie heeft. Zodra een systeem een lange termijn leeft kan het zo zijn dat het systeem anders wordt gebruikt, hier moet de strategie op voorbereid zijn. Een strategie moet moet niet te afhankelijk zijn van het feit dat gebruikers weten hoe het systeem wordt geïmplementeerd. Ook moet de strategie voorbereid zijn op de vooruitgang van hardware.

\subsubsection{Methoden voor schaalbaarheidsborging}
In het onderzoek vertellen Weinstock en Goodenough dat niet echt mogelijk is om te testen of een systeem schaalbaar is. Wel zijn er methoden om de schaalbaarheid te waarborgen, zoals:
\begin{itemize}
	\item Onderzoek de \enquote{performance curves} en karakteriseer deze met een Big O notatie.
	\item Identificeer mechanismen om knelpunten aan het licht brengen of waar aannames van het schaalbaarheids- ontwerp beginnen te worden geschonden.
	\item Voer een SWOT analyse uit op de schaalbaarheids-strategie.
	\begin{itemize}
		\item \textbf{S}trengths (de soorten groei waar de strategie voor ontworpen is)
		\item \textbf{W}eaknesses (de soorten groei waar de strategie niet voor ontworpen is)
		\item \textbf{O}pportunities (mogelijke veranderingen in werklast of technologie die de strategie goed zou kunnen benutten)
		\item \textbf{T}hreats (mogelijke veranderingen in de werklast of technologie die de strategie in twijfel zouden kunnen trekken)
	\end{itemize}
\end{itemize}

Schalen kan op twee verschillende manieren, namelijk horizontaal en verticaal. Horizontaal wilt zeggen dat er geschaald wordt door meerdere machines toe te voegen, terwijl verticaal schalen betekent dat er  meer rekenkracht (als bijvoorbeeld een betere CPU of meer RAM) wordt toegevoegd aan een machine.

\subsection{Prestaties versus schaalbaarheid}

\subsection{Kosten versus schaalbaarheid}

\section{Onderhoudbaarheid}
Grubb, P. en Takang, A. A. definiëren in hun boek \enquote{Software Maintenance: Concepts And Practice} onderhoudbaarheid als \enquote{The discipline concerned with changes related to a software system after delivery} \parencite{MaintenanceConcepts}. In 1993 heeft IEEE een \enquote{Standard Glossary of Software Engineering Terminology} opgesteld. Deze begrippenlijst definieert onderhoudbaarheid als \enquote{the ease with which a software system or component can be modified to correct faults, improve performance or other attributes, or adapt to a changed environment} \parencite{SENTerminology}. Deze twee definities komen uiteindelijk op hetzelfde neer. Grubb en Takang noemen het in de context van een discipline, terwijl IEEE het als een kwaliteitseigenschap definieert. Ook specificeren Grubb en Takang het feit dat het alleen ná het opleveren gebeurt. In dit onderzoek wordt uitgegaan van de definitie van IEEE.

Grubb en Takang noemen ook een aantal redenen waarom software moet worden onderhouden:
\begin{itemize}
	\item Ondersteuning van verplichtte upgrades
	\item Ondersteuning van verzoeken van gebruikers om verbeteringen toe te voegen
	\item Om toekomstige onderhoudswerkzaamheden te vergemakkelijken
\end{itemize}

K.K. Aggarwal, Y. Singh en J.K. Chhabra noemen in hun onderzoek een aantal factoren die van invloed zijn op onderhoudbaarheid van software \parencite{MaintainabilityMeasure}:
\begin{itemize}
	\item Leesbaarheid van de broncode
	\item Kwaliteit van de documentatie
	\item Begrijpelijkheid van software
\end{itemize}

ISO 25010 \parencite{ISO25010} definieert onderhoudbaarheid als \enquote{The degree of effectiveness and efficiency with which a product or system can be modified to improve it, correct it or adapt it to changes in environment, and in requirements} en verdeeld het in een vijftal kwaliteitseigenschappen.
\begin{itemize}
	\item Modulariteit
	\item Herbruikbaarheid
	\item Analyseerbaarheid
	\item Wijzigbaarheid
	\item Testbaarheid
\end{itemize}

Een bekende manier om onderhoudbaarheid te meten is de zogenaamde Maintainability Index (MI). Hier is echter veel kritiek op \parencite{MaintainabilityLiteratureReview, WhyNoMI, WhyNoMI2}.

\section{Infrastructuur}
\enquote{The architecture of a software-intensive system is the structure or structures of the system, which comprise software elements, the externally visible properties of those elements, and the relationships among them.} \parencite{ArchitectureStakeholders}.


\section{Overlapping}
Functional scalability: The ability to enhance the system by adding new functionality without disrupting existing activities.




Ook onderhoudbaar

%\include{Chapters/Chapter4} 
%\include{Chapters/Chapter5} 

%----------------------------------------------------------------------------------------
%	THESIS CONTENT - APPENDICES
%----------------------------------------------------------------------------------------

\appendix % Cue to tell LaTeX that the following "chapters" are Appendices

% Include the appendices of the thesis as separate files from the Appendices folder
% Uncomment the lines as you write the Appendices

\renewcommand{\appendixname}{Bijlage}
% Appendix A

\chapter{Frequently Asked Questions} % Main appendix title

\label{AppendixA} % For referencing this appendix elsewhere, use \ref{AppendixA}

\section{How do I change the colors of links?}

The color of links can be changed to your liking using:

{\small\verb!\hypersetup{urlcolor=red}!}, or

{\small\verb!\hypersetup{citecolor=green}!}, or

{\small\verb!\hypersetup{allcolor=blue}!}.

\noindent If you want to completely hide the links, you can use:

{\small\verb!\hypersetup{allcolors=.}!}, or even better: 

{\small\verb!\hypersetup{hidelinks}!}.

\noindent If you want to have obvious links in the PDF but not the printed text, use:

{\small\verb!\hypersetup{colorlinks=false}!}.

%% Appendix B

\chapter{Tabellen} % Main appendix title

\label{BijlageTabellen} 

\section{Beyond the 12-factor app}

\label{TabelFactors}

In deze tabel worden de 15 factoren behandeld.

\begin{longtable}[c]{c p{3cm} p{3.5cm} p{5.5cm}}
	\toprule
	\textbf{Factor} & \textbf{Naam} & \textbf{Gevolg} & \textbf{Waarom?} \\
	\midrule		
	1 & One codebase, one application & Onderhoudbaarheid (Modulariteit) & Een applicatie is een losstaand component waardoor wijzigingen minimale impact hebben op andere componenten. \\
	2 & API first & Structural scalability & Door de API op de eerste rang te zetten van het development proces wordt de mogelijkheid gecreëerd om met elkaars contracten te communiceren zonder interne ontwikkelingsprocessen te verstoren. Zo kunnen veel nieuwe services gemakkelijker worden toegevoegd. \\
	3 & Dependency management & Onderhoudbaarheid (modulariteit) & Gemakkelijk opzetten van project voor nieuwe ontwikkelaars. \\
	4 & Design, build, release, and run & Onderhoudbaarheid (wijzigbaarheid) & Door duidelijke stadia te definiëren worden wijzigingen aan de applicatie sneller in productie geplaatst. \\
	5 & Configuration, credentials, and code & Structural scalability & Environment variabelen zijn niet in omgevingen maar per deployment opgezet, zo maakt de hoeveelheid omgevingen niet uit. \\
	6 & Logs & Onderhoudbaarheid (analyseerbaarheid) & Door logs naar de \texttt{stdout} te sturen is het gemakkelijker om specifieke fouten te vinden, overzicht te creëren en actief meldingen te versturen naar ontwikkelaars. \\
	7 & Disposability & Load scalability & Door processen gemakkelijk te laten stoppen en starten gaat het schalen een stuk sneller. \\
	8 & Backing services & Onderhoudbaarheid (modulariteit) & Door backing services als \enquote{attached resources} te behandelen maakt het niet uit welke techniek er wordt gebruikt en zijn deze dus los gekoppeld. \\
	9 & Environment parity & Onderhoudbaarheid (wijzigbaarheid) & Er kan een stuk vaker gedeployed worden naar een specifieke omgeving, doordat alle omgevingen zo goed als gelijk aan elkaar zijn. \\
	10 & Administrative processes & Onderhoudbaarheid (analyseerbaarheid) & Door commands in versiebeheer op te slaan is er een duidelijk overzicht en een geschiedenis van alle \enquote{one-off processes} die gebeuren. \\
	11 & Port binding & Onderhoudbaarheid (wijzigbaarheid) & Door HTTP als een service te beschouwen ontstaat er meer controle over lagere levels van de infrastructuur (HTTP \& TCP). \\
	12 & Stateless processes & Load scalability & Mede door de shared-nothing architectuur kan het systeem gemakkelijker schalen. \\
	13 & Concurrency & Load scalability &  Door het horizontaal of verticaal schalen kan de applicatie een groeiende hoeveelheid verkeer beter aan. \\
	14 & Telemetry & Onderhoudbaarheid (Analyseerbaarheid) & Door gegevens van de applicatie in productie goed te kunnen monitoren is op te maken hoe de applicatie zich gedraagt. Zodra er iets fout is kan er meteen op worden gereageerd. \\
	15 & Authentication and authorization & Security & Een cloud-native applicatie moet veilig zijn, aangezien de code over meerdere data centers wordt getransporteerd en door veel verschillende cliënten wordt benaderd. \\
	\bottomrule\\
\end{longtable}

\label{Bijlagek6}
%% Appendix B

\chapter{Feedback} % Main appendix title

\label{Feedback} 

\section{(niet-)functionele schaalbaarheid}

\label{feedbackschaalbaarheid}

Na het promoten van de geschreven blog\footnote{https://developers.nl/blog/69/Defining-software-scalability-using-requirements} is dit de meest populaire blog van Developers.nl in 2019 geworden. Dit heeft het volgende feedbackpuntje opgeleverd: \enquote{Scalability is a two way thing, so adding and removing should be in the definition (and thought lines)}. Waar ik het volledig mee eens ben, en zal verbeteren in de toekomst. Ook is er een leuke discussie uit gekomen. Marlon Etheredge, MSc schreef:
\begin{minted}[bgcolor=codebg, breaklines, breaksymbolleft=]{text}
Hi Kaj,

Mijn vraag behoeft enige introductie.

Ik ben werkzaam in een deelgebied van de informatica/software-engineering waarbij performantie zeer belangrijk is, computer graphics. Onze implementaties dienen zo snel als mogelijk antwoorden te geven op soms complexe berekeningen, doorgaans in (minder dan) millisecondes.

Schaalbaarheid in mijn context staat dan ook voor twee dingen; ten eerste gaat het om het niet schenden van tijdsgrenzen waar wij mee te maken hebben (e.g. één frame dient in 1/50 seconde klaar te zijn) onafhankelijk van de hoeveelheid data die verwerkt moet worden. Ten tweede staat schaalbaarheid voor implementaties die rekbaar zijn op basis van veranderende eisen die aan een systeem worden gesteld.

In deze context gaat het dan niet zo zeer om een veranderend systeem, waarbij bijvoorbeeld functionaliteit wordt toegevoegd
("... must be modified as soon ..." in je eerste definitie), of het systeem verbeterd wordt ("... is able to be improved ..." in je tweede definitie), maar eerder om een kwaliteitskenmerk van een systeem in ogenschouw nemend welke eisen mogelijk in de toekomst aan dit systeem gesteld zullen worden en de hoeveelheid energie die het zal kosten om het systeem te laten aansluiten op deze eisen. In die zin denk ik overigens ook dat dit een interessant onderwerp is, aangezien het ontwerpen van dergelijke systemen fundamenteel is aan de informatica.

Mijn concrete vraag aan jou is als volgt; je schrijft:

"Instead of adding more of the same requirement, non-functional requirements like security or usability are always able to be improved. Therefore, scaling a non-functional requirement is the same as improving it. Setting clear requirements helps proving your solution is scalable non-functionally."

Is verandering (bijvoorbeeld in de vorm van verbetering) noodzakelijk voor schaalbaarheid, of is het mogelijk schaalbaarheid te toetsen los van verandering?
\end{minted}
Mijn antwoord:
\begin{minted}[bgcolor=codebg,breaklines, breaksymbolleft=]{text}
Schaalbaarheid in jouw context sluit goed aan op mijn twee definities: Je noemt "het niet schenden van tijdsgrenzen ... onafhankelijk van de hoeveelheid data"; in dit geval zijn de berekeningen een functionele requirement, en deze moeten voldoende blijven functioneren naarmate het hoeveelheid gebruik toe neemt. In deze context gaat het functioneren dus over het niet schenden van tijdsgrenzen, en de hoeveelheid gebruik over de hoeveelheid data.

De tweede definitie die je noemt (schaalbaarheid voor implementaties die rekbaar zijn op basis van veranderende eisen) omvat in dit geval zowel functionele als niet functionele schaalbaarheid. In mijn definities heb ik het vooral over opschalen, dit is nog een verbeterpunt. Het concept van "veranderende eisen" vind ik een mooie.

Om antwoord te geven op je vraag: Schaalbaarheid is een kwaliteitskenmerk, het daadwerkelijk schalen is een uitoefening van dit kenmerk. Dus, ja, het is mogelijk om schaalbaarheid te toetsen los van daadwerkelijke verandering. Een kwaliteitsanalyse op kenmerken als complexiteit van algoritmes bijvoorbeeld. In de context van functionele schaalbaarheid is dit "to what extent it continues to function properly as the amount of use of the system increases", en van niet-functionele schaalbaarheid "to what extent the quality of that requirement remains acceptable as the use of the system increases".

Ik hoop dat ik hiermee je vraag voldoende heb beantwoord, zo niet hoor ik het graag uiteraard.
\end{minted}



%----------------------------------------------------------------------------------------
%	BIBLIOGRAPHY
%----------------------------------------------------------------------------------------

\renewcommand{\bibname}{Literatuurlijst}
\printbibliography[heading=bibintoc]

%----------------------------------------------------------------------------------------

\end{document}  
